% Emacs settings: -*-mode: latex; TeX-master: "Component_manual"; -*-

\chapter{The Union framework: sample environments and multiple scattering}
\label{c:union}
\index{Library!Components!union}
\index{Union|textbf}

The Union framework, originally contributed by Mads Bertelsen for \MCS
and ported to \MCX, is a set of components that together allow the
description of complex, possibly nested or overlapping, sample
environments with multiple scattering, by separating the description of
\emph{geometry} (where matter is) from \emph{physical process} (what
happens to a photon there, e.g. Compton or Rayleigh scattering,
photoelectric absorption, or powder diffraction).

A Union simulation is assembled in four steps, each using one or more of
the component classes documented in this chapter:
\begin{enumerate}
\item One or more \emph{process} components (e.g. \texttt{Incoherent\_process},
  \texttt{Compton\_xrl\_process}, \texttt{Rayleigh\_xrl\_process},
  \texttt{Powder\_process}) each describe one interaction process.
\item These are gathered into named \emph{materials} using
  \texttt{Union\_make\_material}, which may combine several processes into
  one material definition.
\item \emph{Geometry} components (\texttt{Union\_box}, \texttt{Union\_cylinder},
  \texttt{Union\_sphere}, \texttt{Union\_cone}) place volumes of a given
  material in the instrument, and may be nested or intersected to build up
  arbitrarily complex sample environments.
\item A \texttt{Union\_master} component, placed after all of the geometries
  it should simulate, performs the actual ray-tracing (including multiple
  scattering) through every geometry defined since the previous master (or
  since \texttt{Union\_init}).
\end{enumerate}

\section{Union framework core: initialization, master, and termination}
\section{The \texttt{Union\_init} McStas Component}
Initialize component that needs to be place before any Union component

\subsection*{Identification}
\begin{itemize}
  \item \textbf{Site:} 
  \item \textbf{Author:} Mads Bertelsen
  \item \textbf{Origin:} ESS DMSC
  \item \textbf{Date:} 20.08.15
\end{itemize}

\subsection*{Description}
\begin{lstlisting}
Part of the Union components, a set of components that work together and thus
sperates geometry and physics within McStas.
The use of this component requires other components to be used.

1) One specifies a number of processes using process components
2) These are gathered into material definitions using this component
3) Geometries are placed using Union_box/cylinder/sphere, assigned a material
4) A Union_master component placed after all of the above

Only in step 4 will any simulation happen, and per default all geometries
defined before the master, but after the previous will be simulated here.

There is a dedicated manual available for the Union_components

Algorithm:
Described elsewhere
\end{lstlisting}

\subsection*{Input parameters}
Parameters in \textbf{boldface} are required; the others are optional.

\begin{longtable}{p{0.22\textwidth}p{0.12\textwidth}p{0.46\textwidth}p{0.14\textwidth}}
\toprule
\textbf{Name} & \textbf{Unit} & \textbf{Description} & \textbf{Default} \\
\midrule
\endhead
\bottomrule
\end{longtable}

\subsection*{Links}
\begin{itemize}
  \item \href{run:/home/willend/willend-McCode/mcstas-comps/union/Union_init.comp}{Source code} for \texttt{Union\_init.comp}.
\end{itemize}
\IfFileExists{Union_init_static.tex}{\input{Union_init_static.tex}}{}
\section{The \texttt{Union\_master} McStas Component}
Master component for Union components that performs the overall simulation

\subsection*{Identification}
\begin{itemize}
  \item \textbf{Site:} 
  \item \textbf{Author:} Mads Bertelsen
  \item \textbf{Origin:} University of Copenhagen
  \item \textbf{Date:} 20.08.15
\end{itemize}

\subsection*{Description}
\begin{lstlisting}
Part of the Union components, a set of components that work together and thus
sperates geometry and physics within McStas.
The use of this component requires other components to be used.

1) One specifies a number of processes using process components
2) These are gathered into material definitions using Union_make_material
3) Geometries are placed using Union_box/cylinder/sphere, assigned a material
4) This master component placed after all of the above

Only in step 4 will any simulation happen, and per default all geometries
defined before this master, but after the previous will be simulated here.

There is a dedicated manual available for the Union_components

Algorithm:
Described elsewhere
\end{lstlisting}

\subsection*{Input parameters}
Parameters in \textbf{boldface} are required; the others are optional.

\begin{longtable}{p{0.22\textwidth}p{0.12\textwidth}p{0.46\textwidth}p{0.14\textwidth}}
\toprule
\textbf{Name} & \textbf{Unit} & \textbf{Description} & \textbf{Default} \\
\midrule
\endhead
enable\_refraction & 0/1 & Toggles refraction effects, simulated for all materials that have refraction parameters set & 1 \\
enable\_reflection & 0/1 & Toggles reflection effects, simulated for all materials that have refraction parameters set & 1 \\
verbal & 0/1 & Toogles printing information loaded during initialize & 0 \\
list\_verbal & 0/1 & Toogles information of all internal lists in intersection network & 0 \\
finally\_verbal & 0/1 & Toogles information about cleanup performed in finally section & 0 \\
allow\_inside\_start & 0/1 & Set to 1 if rays are expected to start inside a volume in this master & 0 \\
enable\_tagging & 0/1 & Enable tagging of ray history (geometry, scattering process) & 0 \\
history\_limit & 1 & Limit the number of unique histories that are saved & 300000 \\
enable\_conditionals & 0/1 & Use conditionals with this master & 1 \\
inherit\_number\_of\_scattering\_events & 0/1 & Inherit the number of scattering events from last master & 0 \\
weight\_ratio\_limit & 1 & Limit on the ratio between initial weight and current, ray absorbed if ratio becomes lower than limit & 1e-90 \\
init & string & Name of Union\_init component (typically "init", default) & "init" \\
\bottomrule
\end{longtable}

\subsection*{Links}
\begin{itemize}
  \item \href{run:/home/willend/willend-McCode/mcstas-comps/union/Union_master.comp}{Source code} for \texttt{Union\_master.comp}.
\end{itemize}
\IfFileExists{Union_master_static.tex}{\input{Union_master_static.tex}}{}
\section{The \texttt{Union\_stop} McStas Component}
Stop component that must be placed after all Union components for instrument to compile correctly

\subsection*{Identification}
\begin{itemize}
  \item \textbf{Site:} 
  \item \textbf{Author:} Mads Bertelsen
  \item \textbf{Origin:} University of Copenhagen
  \item \textbf{Date:} 20.08.15
\end{itemize}

\subsection*{Description}
\begin{lstlisting}
Part of the Union components, a set of components that work together and thus
sperates geometry and physics within McStas.
The use of this component requires other components to be used.

1) One specifies a number of processes using process components
2) These are gathered into material definitions using this component
3) Geometries are placed using Union_box/cylinder/sphere, assigned a material
4) A Union_master component placed after all of the above

Only in step 4 will any simulation happen, and per default all geometries
defined before the master, but after the previous will be simulated here.

There is a dedicated manual available for the Union_components

Algorithm:
Described elsewhere
\end{lstlisting}

\subsection*{Input parameters}
Parameters in \textbf{boldface} are required; the others are optional.

\begin{longtable}{p{0.22\textwidth}p{0.12\textwidth}p{0.46\textwidth}p{0.14\textwidth}}
\toprule
\textbf{Name} & \textbf{Unit} & \textbf{Description} & \textbf{Default} \\
\midrule
\endhead
\bottomrule
\end{longtable}

\subsection*{Links}
\begin{itemize}
  \item \href{run:/home/willend/willend-McCode/mcstas-comps/union/Union_stop.comp}{Source code} for \texttt{Union\_stop.comp}.
\end{itemize}
\IfFileExists{Union_stop_static.tex}{\input{Union_stop_static.tex}}{}

\section{Union geometry components}
\section{The \texttt{Union\_box} McStas Component}
A box geometry component for the Union components

\subsection*{Identification}
\begin{itemize}
  \item \textbf{Site:} 
  \item \textbf{Author:} Mads Bertelsen
  \item \textbf{Origin:} University of Copenhagen
  \item \textbf{Date:} 20.08.15
\end{itemize}

\subsection*{Description}
\begin{lstlisting}
Part of the Union components, a set of components that work together and thus
sperates geometry and physics within McStas.
The use of this component requires other components to be used.

1) One specifies a number of processes using process components
2) These are gathered into material definitions using Union_make_material
3) Geometries are placed using Union_box/cylinder/sphere, assigned a material
4) A Union_master component placed after all of the above

Only in step 4 will any simulation happen, and per default all geometries
defined before this master, but after the previous will be simulated here.

There is a dedicated manual available for the Union components

The position of this component is the center of the box, zdepth/2 in each direction.

It is allowed to overlap components, but it is not allowed to have two
parallel planes that coincides. This will crash the code on run time.


Algorithm:
Described elsewhere
\end{lstlisting}

\subsection*{Input parameters}
Parameters in \textbf{boldface} are required; the others are optional.

\begin{longtable}{p{0.22\textwidth}p{0.12\textwidth}p{0.46\textwidth}p{0.14\textwidth}}
\toprule
\textbf{Name} & \textbf{Unit} & \textbf{Description} & \textbf{Default} \\
\midrule
\endhead
material\_string & string & Material name of this volume, defined using Union\_make\_material & 0 \\
\textbf{priority} & 1 & Priotiry of the volume (can not be the same as another volume) A high priority is on top of low. &  \\
\textbf{xwidth} & m & Width of the box volume &  \\
\textbf{yheight} & m & Height of the box volume &  \\
\textbf{zdepth} & m & Depth of the box volume &  \\
xwidth2 & m & Optional different width at the +z box face & -1 \\
yheight2 & m & Optional different height at the +z box face & -1 \\
visualize & 1 & Set to 0 if you wish to hide this geometry in mcdisplay & 1 \\
target\_index & string & Focuses on component a component this many steps further in the component sequence & 0 \\
target\_x & m & X position of target to focus at & 0 \\
target\_y & m & Y position of target to focus at & 0 \\
target\_z & m & Z position of target to focus at & 0 \\
focus\_aw & deg & Horiz. angular dimension of a rectangular area & 0 \\
focus\_ah & deg & Vert. angular dimension of a rectangular area & 0 \\
focus\_xw & m & Horiz. dimension of a rectangular area & 0 \\
focus\_xh & m & Vert. dimension of a rectangular area & 0 \\
focus\_r & m & Focusing on circle with this radius & 0 \\
plus\_z\_surface & string & Comma seperated string of Union surface definitions defined prior to this component, applied to +z face of box & 0 \\
minus\_z\_surface & string & Comma seperated string of Union surface definitions defined prior to this component, applied to -z face of box & 0 \\
plus\_x\_surface & string & Comma seperated string of Union surface definitions defined prior to this component, applied to +x face of box & 0 \\
minus\_x\_surface & string & Comma seperated string of Union surface definitions defined prior to this component, applied to -x face of box & 0 \\
plus\_y\_surface & string & Comma seperated string of Union surface definitions defined prior to this component, applied to +y face of box & 0 \\
minus\_y\_surface & string & Comma seperated string of Union surface definitions defined prior to this component, applied to -z face of box & 0 \\
all\_face\_surface & string & Comma seperated string of Union surface definitions defined prior to this component, applied to all faces above & 0 \\
cut\_surface & string & Comma seperated string of Union surface definitions defined prior to this component, applied to all internal cuts & 0 \\
p\_interact & 1 & Probability to interact with this geometry [0-1] & 0 \\
mask\_string & string & Comma seperated list of geometry names which this geometry should mask & 0 \\
mask\_setting & string & "All" or "Any", should the masked volume be simulated when the ray is in just one mask, or all. & 0 \\
number\_of\_activations & 1 & Number of subsequent Union\_master components that will simulate this geometry & 1 \\
init & string & name of Union\_init component (typically "init", default) & "init" \\
\bottomrule
\end{longtable}

\subsection*{Links}
\begin{itemize}
  \item \href{run:/home/willend/willend-McCode/mcstas-comps/union/Union_box.comp}{Source code} for \texttt{Union\_box.comp}.
\end{itemize}
\IfFileExists{Union_box_static.tex}{\input{Union_box_static.tex}}{}
\section{The \texttt{Union\_cylinder} McStas Component}
Cylinder geometry component for Union components

\subsection*{Identification}
\begin{itemize}
  \item \textbf{Site:} 
  \item \textbf{Author:} Mads Bertelsen
  \item \textbf{Origin:} University of Copenhagen
  \item \textbf{Date:} 20.08.15
\end{itemize}

\subsection*{Description}
\begin{lstlisting}
Part of the Union components, a set of components that work together and thus
sperates geometry and physics within McStas.
The use of this component requires other components to be used.

1) One specifies a number of processes using process components
2) These are gathered into material definitions using Union_make_material
3) Geometries are placed using Union_box/cylinder/sphere, assigned a material
4) A Union_master component placed after all of the above

Only in step 4 will any simulation happen, and per default all geometries
defined before this master, but after the previous will be simulated here.

There is a dedicated manual available for the Union_components

The position of this component is the center of the cylinder, and it thus
extends yheight/2 up and down along y axis.

It is allowed to overlap components, but it is not allowed to have two
parallel planes that coincides. This will crash the code on run time.
\end{lstlisting}

\subsection*{Input parameters}
Parameters in \textbf{boldface} are required; the others are optional.

\begin{longtable}{p{0.22\textwidth}p{0.12\textwidth}p{0.46\textwidth}p{0.14\textwidth}}
\toprule
\textbf{Name} & \textbf{Unit} & \textbf{Description} & \textbf{Default} \\
\midrule
\endhead
material\_string & string & Material name of this volume, defined using Union\_make\_material & 0 \\
\textbf{priority} & 1 & Priotiry of the volume (can not be the same as another volume) A high priority is on top of low. &  \\
\textbf{radius} & m & Outer radius volume in (x,z) plane &  \\
\textbf{yheight} & m & Cylinder height in (y) direction &  \\
visualize & 1 & Set to 0 if you wish to hide this geometry in mcdisplay & 1 \\
target\_index & 1 & Focuses on component a component this many steps further in the component sequence & 0 \\
target\_x & m & X Position of target to focus at & 0 \\
target\_y & m & Y Position of target to focus at & 0 \\
target\_z & m & Z Position of target to focus at & 0 \\
focus\_aw & deg & Horiz. angular dimension of a rectangular area & 0 \\
focus\_ah & deg & Vert. angular dimension of a rectangular area & 0 \\
focus\_xw & m & Horiz. dimension of a rectangular area & 0 \\
focus\_xh & m & Vert. dimension of a rectangular area & 0 \\
focus\_r & m & Focusing on circle with this radius & 0 \\
p\_interact & 1 & Probability to interact with this geometry [0-1] & 0 \\
mask\_string & string & Comma seperated list of geometry names which this geometry should mask & 0 \\
mask\_setting & string & "All" or "Any", should the masked volume be simulated when the ray is in just one mask, or all. & 0 \\
number\_of\_activations & 1 & Number of subsequent Union\_master components that will simulate this geometry & 1 \\
curved\_surface & string & Comma seperated string of Union surface definitions defined prior to this component, applied to curved wall of cylinder & 0 \\
top\_surface & string & Comma seperated string of Union surface definitions defined prior to this component, applied to +y top face of cylinder & 0 \\
bottom\_surface & string & Comma seperated string of Union surface definitions defined prior to this component, applied to -y bottom face of cylinder & 0 \\
all\_face\_surface & string & Comma seperated string of Union surface definitions defined prior to this component, applied to all faces above & 0 \\
cut\_surface & string & Comma seperated string of Union surface definitions defined prior to this component, applied to all internal cuts & 0 \\
init & string & Name of Union\_init component (typically "init", default) & "init" \\
\bottomrule
\end{longtable}

\subsection*{Links}
\begin{itemize}
  \item \href{run:/home/willend/willend-McCode/mcstas-comps/union/Union_cylinder.comp}{Source code} for \texttt{Union\_cylinder.comp}.
\end{itemize}
\IfFileExists{Union_cylinder_static.tex}{\input{Union_cylinder_static.tex}}{}
\section{The \texttt{Union\_sphere} McStas Component}
Sphere geometry component for the Union components

\subsection*{Identification}
\begin{itemize}
  \item \textbf{Site:} 
  \item \textbf{Author:} Mads Bertelsen
  \item \textbf{Origin:} University of Copenhagen
  \item \textbf{Date:} 20.08.15
\end{itemize}

\subsection*{Description}
\begin{lstlisting}
Part of the Union components, a set of components that work together and thus
sperates geometry and physics within McStas.
The use of this component requires other components to be used.

1) One specifies a number of processes using process components
2) These are gathered into material definitions using Union_make_material
3) Geometries are placed using Union_box/cylinder/sphere, assigned a material
4) A Union_master component placed after all of the above

Only in step 4 will any simulation happen, and per default all geometries
defined before this master, but after the previous will be simulated here.

There is a dedicated manual available for the Union components

The position of this component is the center of the sphere

It is allowed to overlap components, but it is not allowed to have two
parallel planes that coincides. This will crash the code on run time.
\end{lstlisting}

\subsection*{Input parameters}
Parameters in \textbf{boldface} are required; the others are optional.

\begin{longtable}{p{0.22\textwidth}p{0.12\textwidth}p{0.46\textwidth}p{0.14\textwidth}}
\toprule
\textbf{Name} & \textbf{Unit} & \textbf{Description} & \textbf{Default} \\
\midrule
\endhead
material\_string & string & material name of this volume, defined using Union\_make\_material & 0 \\
\textbf{priority} & 1 & priotiry of the volume (can not be the same as another volume) A high priority is on top of low. &  \\
\textbf{radius} & m & Radius of sphere &  \\
visualize & 1 & set to 0 if you wish to hide this geometry in mcdisplay & 1 \\
target\_index & 1 & Focuses on component a component this many steps further in the component sequence & 0 \\
target\_x & m & X position of target to focus at & 0 \\
target\_y & m & Y position of target to focus at & 0 \\
target\_z & m & Z position of target to focus at & 0 \\
focus\_aw & deg & horiz. angular dimension of a rectangular area & 0 \\
focus\_ah & deg & vert. angular dimension of a rectangular area & 0 \\
focus\_xw & m & horiz. dimension of a rectangular area & 0 \\
focus\_xh & m & vert. dimension of a rectangular area & 0 \\
focus\_r & m & focusing on circle with this radius & 0 \\
p\_interact & 1 & probability to interact with this geometry [0-1] & 0 \\
mask\_string & string & Comma seperated list of geometry names which this geometry should mask & 0 \\
mask\_setting & string & "All" or "Any", should the masked volume be simulated when the ray is in just one mask, or all. & 0 \\
number\_of\_activations & 1 & Number of subsequent Union\_master components that will simulate this geometry & 1 \\
surface & string & Comma seperated string of Union surface definitions defined prior to this component & 0 \\
cut\_surface & string & Comma seperated string of Union surface definitions defined prior to this component, applied to all internal cuts & 0 \\
init & string & name of Union\_init component (typically "init", default) & "init" \\
\bottomrule
\end{longtable}

\subsection*{Links}
\begin{itemize}
  \item \href{run:/home/willend/willend-McCode/mcstas-comps/union/Union_sphere.comp}{Source code} for \texttt{Union\_sphere.comp}.
\end{itemize}
\IfFileExists{Union_sphere_static.tex}{\input{Union_sphere_static.tex}}{}
\section{The \texttt{Union\_cone} McStas Component}
Cone geometry component for Union components

\subsection*{Identification}
\begin{itemize}
  \item \textbf{Site:} 
  \item \textbf{Author:} Martin Olsen
  \item \textbf{Origin:} University of Copenhagen
  \item \textbf{Date:} 17.09.18
\end{itemize}

\subsection*{Description}
\begin{lstlisting}
Part of the Union components, a set of components that work together and thus
sperates geometry and physics within McStas.
The use of this component requires other components to be used.

1) One specifies a number of processes using process components
2) These are gathered into material definitions using Union_make_material
3) Geometries are placed using Union_box/cylinder/sphere, assigned a material
4) A Union_master component placed after all of the above

Only in step 4 will any simulation happen, and per default all geometries
defined before this master, but after the previous will be simulated here.

There is a dedicated manual available for the Union_components

The position of this component is the center of the cone, and it thus
extends yheight/2 up and down along y axis.

It is allowed to overlap components, but it is not allowed to have two
parallel planes that coincides. This will crash the code on run time.
\end{lstlisting}

\subsection*{Input parameters}
Parameters in \textbf{boldface} are required; the others are optional.

\begin{longtable}{p{0.22\textwidth}p{0.12\textwidth}p{0.46\textwidth}p{0.14\textwidth}}
\toprule
\textbf{Name} & \textbf{Unit} & \textbf{Description} & \textbf{Default} \\
\midrule
\endhead
material\_string & string & Material name of this volume, defined using Union\_make\_material & 0 \\
\textbf{priority} & 1 & Priotiry of the volume (can not be the same as another volume) A high priority is on top of low. &  \\
radius & m & Radius volume in (x,z) plane & 0 \\
radius\_top & m & Top radius volume in (x,z) plane & 0 \\
radius\_bottom & m & Bottom radius volume in (x,z) plane & 0 \\
\textbf{yheight} & m & Cone height in (y) direction &  \\
visualize & 1 & Set to 0 if you wish to hide this geometry in mcdisplay & 1 \\
target\_index & 1 & Focuses on component a component this many steps further in the component sequence & 0 \\
target\_x & m & X position of target to focus at & 0 \\
target\_y & m & Y position of target to focus at & 0 \\
target\_z & m & Z position of target to focus at & 0 \\
focus\_aw & deg & Horiz. angular dimension of a rectangular area & 0 \\
focus\_ah & deg & Vert. angular dimension of a rectangular area & 0 \\
focus\_xw & m & Horiz. dimension of a rectangular area & 0 \\
focus\_xh & m & Vert. dimension of a rectangular area & 0 \\
focus\_r & m & Focusing on circle with this radius & 0 \\
p\_interact & 1 & Probability to interact with this geometry [0-1] & 0 \\
mask\_string & string & Comma seperated list of geometry names which this geometry should mask & 0 \\
mask\_setting & string & "All" or "Any", should the masked volume be simulated when the ray is in just one mask, or all. & 0 \\
number\_of\_activations & 1 & Number of subsequent Union\_master components that will simulate this geometry & 1 \\
curved\_surface & string & Comma seperated string of Union surface definitions defined prior to this component, applied to curved wall of cone & 0 \\
top\_surface & string & Comma seperated string of Union surface definitions defined prior to this component, applied to +y top face of cone & 0 \\
bottom\_surface & string & Comma seperated string of Union surface definitions defined prior to this component, applied to -y bottom face of cone & 0 \\
all\_face\_surface & string & Comma seperated string of Union surface definitions defined prior to this component, applied to all faces above & 0 \\
cut\_surface & string & Comma seperated string of Union surface definitions defined prior to this component, applied to all internal cuts & 0 \\
init & string & Name of Union\_init component (typically "init", default) & "init" \\
\bottomrule
\end{longtable}

\subsection*{Links}
\begin{itemize}
  \item \href{run:/home/willend/willend-McCode/mcstas-comps/union/Union_cone.comp}{Source code} for \texttt{Union\_cone.comp}.
\end{itemize}
\IfFileExists{Union_cone_static.tex}{\input{Union_cone_static.tex}}{}

\section{Union material definition}
\section{The \texttt{Union\_make\_material} McStas Component}
Component that takes a number of Union processes and constructs a Union material

\subsection*{Identification}
\begin{itemize}
  \item \textbf{Site:} 
  \item \textbf{Author:} Mads Bertelsen
  \item \textbf{Origin:} University of Copenhagen
  \item \textbf{Date:} 20.08.15
\end{itemize}

\subsection*{Description}
\begin{lstlisting}
Part of the Union components, a set of components that work together and thus
sperates geometry and physics within McStas.
The use of this component requires other components to be used.

1) One specifies a number of processes using process components
2) These are gathered into material definitions using this component
3) Geometries are placed using Union_box/cylinder/sphere, assigned a material
4) A Union_master component placed after all of the above

Only in step 4 will any simulation happen, and per default all geometries
defined before the master, but after the previous will be simulated here.

There is a dedicated manual available for the Union_components

Algorithm:
Described elsewhere
\end{lstlisting}

\subsection*{Input parameters}
Parameters in \textbf{boldface} are required; the others are optional.

\begin{longtable}{p{0.22\textwidth}p{0.12\textwidth}p{0.46\textwidth}p{0.14\textwidth}}
\toprule
\textbf{Name} & \textbf{Unit} & \textbf{Description} & \textbf{Default} \\
\midrule
\endhead
process\_string & string & Comma seperated names of physical processes & "NULL" \\
\textbf{my\_absorption} & 1/m & Inverse penetration depth from absorption at standard energy &  \\
absorber & 0/1 & Control parameter, if set to 1 the material will have no scattering processes & 0 \\
refraction\_density & g/cm3 & Density of the refracting material. density \textless{} 0 means the material is outside/before the shape. & 0 \\
refraction\_sigma\_coh & barn & Coherent cross section of refracting material. Use negative value to indicate a negative coherent scattering length & 0 \\
refraction\_weight & g/mol & Molar mass of the refracting material & 0 \\
refraction\_SLD & AA\textasciicircum{}-2 & Scattering length density of material (overwrites sigma\_coh, density and weight) & -1500 \\
init & string & Name of Union\_init component (typically "init", default) & "init" \\
\bottomrule
\end{longtable}

\subsection*{Links}
\begin{itemize}
  \item \href{run:/home/willend/willend-McCode/mcstas-comps/union/Union_make_material.comp}{Source code} for \texttt{Union\_make\_material.comp}.
\end{itemize}
\IfFileExists{Union_make_material_static.tex}{\input{Union_make_material_static.tex}}{}

\section{Union scattering/interaction process components}
\input{union/Incoherent_process}
\input{union/Compton_xrl_process}
\input{union/Rayleigh_xrl_process}
\input{union/KN_xrl_process}
\section{The \texttt{Powder\_process} McStas Component}
Port of the PowderN process to the Union components

\subsection*{Identification}
\begin{itemize}
  \item \textbf{Site:} 
  \item \textbf{Author:} Mads Bertelsen
  \item \textbf{Origin:} University of Copenhagen
  \item \textbf{Date:} 20.08.15
\end{itemize}

\subsection*{Description}
\begin{lstlisting}
This Union_process is based on the PowderN.comp component originally written
by P. Willendrup, L. Chapon, K. Lefmann, A.B.Abrahamsen, N.B.Christensen,
E.M.Lauridsen.

Part of the Union components, a set of components that work together and thus
sperates geometry and physics within McStas.
The use of this component requires other components to be used.

1) One specifies a number of processes using process components like this one
2) These are gathered into material definitions using Union_make_material
3) Geometries are placed using Union_box / Union_cylinder, assigned a material
4) A Union_master component placed after all of the above

Only in step 4 will any simulation happen, and per default all geometries
defined before the master, but after the previous will be simulated here.

There is a dedicated manual available for the Union_components
Algorithm:
Described elsewhere
\end{lstlisting}

\subsection*{Input parameters}
Parameters in \textbf{boldface} are required; the others are optional.

\begin{longtable}{p{0.22\textwidth}p{0.12\textwidth}p{0.46\textwidth}p{0.14\textwidth}}
\toprule
\textbf{Name} & \textbf{Unit} & \textbf{Description} & \textbf{Default} \\
\midrule
\endhead
reflections & string & Input file for reflections. No scattering if NULL or "" [string] & "NULL" \\
packing\_factor & 1 & How dense is the material compared to optimal 0-1 & 1 \\
Vc & AA\textasciicircum{}3 & Volume of unit cell=nb atoms per cell/density of atoms. & 0 \\
delta\_d\_d & 0/1 & Global relative delta\_d\_d/d broadening when the 'w' column is not available. Use 0 if ideal. & 0 \\
DW & 1 & Global Debye-Waller factor when the 'DW' column is not available. Use 1 if included in F2 & 0 \\
nb\_atoms & 1 & Number of sub-unit per unit cell, that is ratio of sigma for chemical formula to sigma per unit cell & 1 \\
d\_phi & deg & Angle corresponding to the vertical angular range to focus to, e.g. detector height. 0 for no focusing. & 0 \\
density & g/cm\textasciicircum{}3 & Density of material. rho=density/weight/1e24*N\_A. & 0 \\
weight & g/mol & Atomic/molecular weight of material. & 0 \\
barns & 1 & Flag to indicate if |F|\textasciicircum{}2 from 'reflections' is in barns or fm\textasciicircum{}2 (barns=1 for laz, barns=0 for lau type files). & 1 \\
Strain & ppm & Global relative delta\_d\_d/d shift when the 'Strain' column is not available. Use 0 if ideal. & 0 \\
interact\_fraction & 1 & How large a part of the scattering events should use this process 0-1 (sum of all processes in material = 1) & -1 \\
format & no quotes & Name of the format, or list of column indexes (see Description). & \{0, 0, 0, 0, 0, 0, 0, 0, 0\} \\
\bottomrule
\end{longtable}

\subsection*{Links}
\begin{itemize}
  \item \href{run:/home/willend/willend-McCode/mcstas-comps/union/Powder_process.comp}{Source code} for \texttt{Powder\_process.comp}.
\end{itemize}
\IfFileExists{Powder_process_static.tex}{\input{Powder_process_static.tex}}{}
\section{The \texttt{Template\_process} McStas Component}
A template process for building Union processes

\subsection*{Identification}
\begin{itemize}
  \item \textbf{Site:} 
  \item \textbf{Author:} Mads Bertelsen
  \item \textbf{Origin:} University of Copenhagen
  \item \textbf{Date:} 20.08.15
\end{itemize}

\subsection*{Description}
\begin{lstlisting}
This is a template for a new contributor to create their own physical process.
The comments in this file are meant to teach the user about creating their own
process file, rather than explaining this one. For comments on how this code works,
look in the Incoherent_process.comp.

Part of the Union components, a set of components that work together and thus
sperates geometry and physics within McStas.
The use of this component requires other components to be used.

1) One specifies a number of processes using process components like this one
2) These are gathered into material definitions using Union_make_material
3) Geometries are placed using Union_box / Union_cylinder, assigned a material
4) A Union_master component placed after all of the above

Only in step 4 will any simulation happen, and per default all geometries
defined before the master, but after the previous will be simulated here.

There is a dedicated manual available for the Union_components


Algorithm:
Described elsewhere
\end{lstlisting}

\subsection*{Input parameters}
Parameters in \textbf{boldface} are required; the others are optional.

\begin{longtable}{p{0.22\textwidth}p{0.12\textwidth}p{0.46\textwidth}p{0.14\textwidth}}
\toprule
\textbf{Name} & \textbf{Unit} & \textbf{Description} & \textbf{Default} \\
\midrule
\endhead
sigma & barns & Incoherent scattering cross section & 5.08 \\
packing\_factor & 1 & How dense is the material compared to optimal 0-1 & 1 \\
unit\_cell\_volume & AA\textasciicircum{}3 & Unit\_cell\_volume & 13.8 \\
interact\_fraction & 1 & How large a part of the scattering events should use this process 0-1 (sum of all processes in material = 1) & -1 \\
init & string & Name of Union\_init component (typically "init", default) & "init" \\
\bottomrule
\end{longtable}

\subsection*{Links}
\begin{itemize}
  \item \href{run:/home/willend/willend-McCode/mcstas-comps/union/Template_process.comp}{Source code} for \texttt{Template\_process.comp}.
\end{itemize}
\IfFileExists{Template_process_static.tex}{\input{Template_process_static.tex}}{}
